\documentclass[11pt,a4paper]{article}
\usepackage[T1]{fontenc}
\usepackage{lmodern}
\usepackage[latin1]{inputenc}
\usepackage{ngerman}
\usepackage{a4wide}
\usepackage[dvips]{graphicx}
\usepackage{float}

\usepackage[
pdfauthor={},
pdftitle={},
pdfcreator={pdftex},
]{hyperref}

\usepackage{sectsty}
\allsectionsfont{\sffamily}

\usepackage{fancyheadings} 
\pagestyle{fancy} 
\lhead{\textsf{\textbf{ACE} \\ \small{a collaborative editor}}}
\chead{}
\rhead{
\parbox[c]{3cm}{\includegraphics[height=0.875cm,width=3cm]{../../images/logo_BFH.eps}}
\parbox[c]{2.2cm}
{\tiny{\textsf{Berner Fachhochschule \\
Hochschule f�r \\
Technik und Informatik}}}}
\lfoot{}
\cfoot{\textsf{\thepage}}
\rfoot{}
\setlength{\headrulewidth}{0.6pt}
\setlength{\footrulewidth}{0.6pt}
\setlength{\topmargin}{-50pt}
\addtolength{\headheight}{50pt}

\usepackage{colortbl}

\newcommand{\headercol}[2]{\multicolumn{1}{|>{\bfseries\columncolor[gray]{0.82}}p{#1}|}{\textsf{#2}}}
\newcommand{\ace}[0]{\emph{ACE }}



\begin{document}

\begin{titlepage}
\thispagestyle{empty}
  \includegraphics[height=1.5in]{../images/pix.eps}

  \begin{center}

    {\fontsize{40}{45} \textbf{\textsf{ACE}}} \\
    \textsf{a collaborative editor} \\
        
    \vspace{36pt}
        
    {\huge{\textbf{\textsf{}}}} \\

    \vspace{36pt}

	\textsf{Berne University of Applied Sciences} \\
    \textsf{School of Engineering and Information Technology} \\
    
  \end{center}

  \vfill
  
  \begin{tabular}{ll}
   \hline

   \\

   \multicolumn{1}{>{\bfseries}p{1.5in}}{\textsf{Date:}} &
   \multicolumn{1}{>{}p{4.3in}}{\textsf{08.11.2005}}          \\
   
   \\
   
   \multicolumn{1}{>{\bfseries}p{1.5in}}{\textsf{Version:}}     &   
   \multicolumn{1}{>{}p{4.3in}}{\textsf{0.1}}                 \\

   \\
   
   \multicolumn{1}{>{\bfseries}p{1.5in}}{\textsf{Projectteam:}}                 &
   \multicolumn{1}{>{}p{4.3in}}{\textsf{Mark Bigler (biglm2@hta-bi.bfh.ch)}}  \\
   \multicolumn{1}{>{\bfseries}p{1.5in}}{}                                      &
   \multicolumn{1}{>{}p{4.3in}}{\textsf{Simon Raess (rasss@hta-bi.bfh.ch)}}    \\
   \multicolumn{1}{>{\bfseries}p{1.5in}}{}                                      &
   \multicolumn{1}{>{}p{4.3in}}{\textsf{Lukas Zbinden (zbinl@hta-bi.bfh.ch)}} \\   
   
   \\
   
   \multicolumn{1}{>{\bfseries}p{1.5in}}{\textsf{Receivers:}}                       &
   \multicolumn{1}{>{}p{4.3in}}{\textsf{Jean-Paul Dubois (doj@hta-bi.bfh.ch)}}       \\
   \multicolumn{1}{>{\bfseries}p{1.5in}}{}                                          &
   \multicolumn{1}{>{}p{4.3in}}{\textsf{Claude Fuhrer (frc@hta-bi.bfh.ch)}}       \\

   \\
   
   \multicolumn{1}{>{\bfseries}p{1.5in}}{\textsf{Location:}}               &   
   \multicolumn{1}{>{}p{4.3in}}{\textsf{Subversion Repository}} \\

   \\  
   
   \hline
  \end{tabular}

\end{titlepage}

\newpage

\tableofcontents
\newpage

\section*{Versionskontrolle}

\begin{table}[!h]
 \begin{tabular}{|l|l|l|l|}
  \hline
  \headercol{0.6in}{Version}         & 
  \headercol{0.8in}{Datum}           &
  \headercol{1.2in}{Verantwortlich}  & 
  \headercol{2.8in}{Bemerkungen}     \\
  \hline
  0.1         & 15.03.2005  & rasss           &  Erste Version \\
  \hline
  0.2         & 22.03.2005  & rasss zbinl     &  �berarbeitung \\
  \hline
  0.3         & 30.03.2005  & zbinl           &  Projektspezifisches Vorgehensmodell \\
  \hline
  0.4         & 30.03.2005  & rasss zbinl     &  �berarbeitung Punkte 2, 3 und 4 \\
  \hline
  0.5         & 05.04.2005  & zbinl           &  �berarbeitung Appendix A.1 \\
  \hline
  0.6         & 06.04.2005  & zbinl           &  Anpassungen Punkte 2, 3 und 4 \\
  \hline
  0.7         & 06.04.2005  & Projektteam     &  Review \\
  \hline
  0.8         & 13.04.2005  & Projektteam     &  Letzte Anpassungen \\
  \hline
  1.0         & 13.04.2005  & Projektteam     &  Freigabe \\
  \hline
 \end{tabular}
 \caption{Versionskontrolle}
 \label{Versionskontrolle}
\end{table}

\begin{table}[!h]
 \begin{tabular}{|l|l|l|l|l|}
  \hline
  \headercol{0.9in}{}            & 
  \headercol{0.9in}{Stelle}      & 
  \headercol{0.8in}{Datum}       & 
  \headercol{0.6in}{Visum}       & 
  \headercol{2.0in}{Bemerkungen} \\
  \hline
  \textbf{Freigegeben}   & Projektteam &       &       &             \\
  \hline
  \textbf{Genehmigt}     &             &       &       &             \\
  \hline
 \end{tabular}
 \caption{Pr�fung/Genehmigung}
 \label{Pr�fung/Genehmigung}
\end{table}

\newpage



\section{Einleitung}
\subsection{Zweck des Dokumentes}
Der Erfahrungsbericht dient der Auswertung der Semesterarbeit im Sommersemester vom Februar bis Juni 2005. Die Ergebnisse werden mit der Planung gem�ss Projektplan verglichen. In einem zweiten Teil werden pers�nliche Erfahrungen und Erkenntnisse festgehalten. Sie liefern wertvolle Informationen, welche im nachfolgenden Diplomprojekt genutzt werden k�nnen, um positive Aspekte zu �bernehmen und negative Vorf�lle m�glichst zu vermeiden.

\subsection{Referenzierte Dokumente}
Die referenzierten Dokumente sind auf der Projekt-Webseite abrufbar: \href{http://ace.iserver.ch/}{http://ace.iserver.ch/}
\begin{itemize}
 \item Projektplan (Version 1.0)
\end{itemize}



\section{Ausgangssituation}
Das Editieren von Dokumenten in einer Gruppe kann eine grosse Herausforderung sein. Versionierungs Systeme wie Subversion oder CVS helfen einer Gruppe eine konsistente Kopie des Dokumentes zu haben, aber erlauben keine Zusammenarbeit in Echtzeit. Das kollaborative Bearbeiten des selben Dokumentes in Echtzeit in einer Gruppe ist heute kaum bekannt, obschon die Idee schon lange erforscht wird. Interessanterweise gibt es aber kaum kommerzielle Editoren, wobei uns SubEthaEdit f�r Mac OS X als einzige Anwendung bekannt ist.

Wir sind der �berzeugung, dass Anwendungen, die es erm�glichen ein Dokument in Echtzeit in einer Gruppe zu bearbeiten ein grosses Potenzial haben. Unser Ziel ist es, den ersten platformunabh�ngigen und voll funktionsf�higen kollaborativen Text-Editor zu erstellen. 

In der Semesterarbeit wollten wir die Grundlage legen f�r die Implementation des Texteditors. Das beinhaltete vorallem einen funktionsf�higen Concurrency Control Algorithmus zu implementieren.



\section{Soll-Ist Vergleich}
Im folgenden wird aufgezeigt, wo es im Projekt zu Differenzen zwischen der Planung und dem tats�chlichen Ablauf des Projektes gekommen ist.

\subsection{Resultate}
W�hrend der Planung wurde festgelegt, welche Resultate w�hrend der Diplomarbeit 
erreicht werden sollen.

\begin{table}[H]
 \begin{tabular}{|l|l|l|}
  \hline
   \headercol{1.4in}{Ergebnisse} &
   \headercol{2in}{Soll} &
   \headercol{2.2in}{Ist} \\
  \hline
   \contentcol{1.4in}{Projektantrag} &
   \contentcol{2in}{Genehmigung Projektantrag} &
   \contentcol{2.2in}{Projektantrag wurde erstellt und genehmigt} \\
  \hline
   \contentcol{1.4in}{Projekthandbuch} &
   \contentcol{2in}{} &
   \contentcol{2.2in}{Projekthandbuch erstellt} \\
  \hline
   \contentcol{1.4in}{Projektplan} &
   \contentcol{2in}{} &
   \contentcol{2.2in}{Projektplan erstellt} \\
  \hline
   \contentcol{1.4in}{Pflichtenheft} &
   \contentcol{2in}{} &
   \contentcol{2.2in}{Pflichtenheft erstellt} \\
  \hline
  \hline
   \contentcol{1.4in}{Bericht Evaluation Algorithmus} &
   \contentcol{2in}{erstellt} &
   \contentcol{2.2in}{abgeschlossen} \\
  \hline
   \contentcol{1.4in}{Algorithmus} &
   \contentcol{2in}{erstellt und dokumentiert} &
   \contentcol{2.2in}{Der Algorithmus ist fertig implementiert. Wir hatten 
   sogar noch die M�glichkeit Undo einzubauen.} \\
  \hline
   \contentcol{1.4in}{Testframework} &
   \contentcol{2in}{erstellt und dokumentiert} &
   \contentcol{2.2in}{Das Testframework ist erstellt und dokumentiert. 
   Eine Anleitung zum Erstellen von Testf�llen ist vorhanden.} \\
  \hline
  \hline
   \contentcol{1.4in}{Bericht Evaluation GUI} &
   \contentcol{2in}{erstellt} &
   \contentcol{2.2in}{erstellt} \\
  \hline
   \contentcol{1.4in}{Bericht Evaluation Netzwerk} &
   \contentcol{2in}{erstellt} &
   \contentcol{2.2in}{erstellt} \\
  \hline
 \end{tabular}
 \caption{Resultate}
 \label{table:resultate}
\end{table}
 

\subsection{Termine}
\begin{table}[H]
 \begin{tabular}{|l|l|l|l}
  \hline
   \headercol{2in}{Ergebnis} &
   \headercol{1in}{Soll} &
   \headercol{1in}{Ist} &
   \headercol{1in}{Status} \\
  \hline
   \contentcol{2in}{Projektantrag} &
   \contentcol{1in}{11.03.2005} &
   \contentcol{1in}{11.03.2005} &
   \contentcol{1in}{abgeschlossen} \\
  \hline
   \contentcol{2in}{Pflichtenheft} &
   \contentcol{1in}{05.04.2005} &
   \contentcol{1in}{05.04.2005} &
   \contentcol{1in}{abgeschlossen} \\
  \hline
   \contentcol{2in}{Projekthandbuch} &
   \contentcol{1in}{15.04.2005} &
   \contentcol{1in}{15.04.2005} &
   \contentcol{1in}{abgeschlossen} \\ 
  \hline
   \contentcol{2in}{Projektplan} &
   \contentcol{1in}{15.04.2005} &
   \contentcol{1in}{15.04.2005} &
   \contentcol{1in}{abgeschlossen} \\
  \hline
  \hline
   \contentcol{2in}{Testframework} &
   \contentcol{1in}{23.05.2005} &
   \contentcol{1in}{23.05.2005} &
   \contentcol{1in}{abgeschlossen} \\
  \hline
   \contentcol{2in}{Algorithmus} &
   \contentcol{1in}{10.06.2005} &
   \contentcol{1in}{???} &
   \contentcol{1in}{???} \\
  \hline
  \hline
   \contentcol{2in}{Bericht GUI} &
   \contentcol{1in}{22.04.2005} &
   \contentcol{1in}{???} &
   \contentcol{1in}{???} \\
  \hline
   \contentcol{2in}{Bericht Netzwerk} &
   \contentcol{1in}{10.06.2005} &
   \contentcol{1in}{06.06.2005} &
   \contentcol{1in}{abgeschlossen} \\
  \hline
 \end{tabular}
 \caption{Termine}
 \label{table:termine}
\end{table}


\subsection{Zeitaufwand}
\begin{table}[H]
 \begin{tabular}{|l|l|l|}
  \headercol{2in}{Ergebnis} &
  \headercol{1in}{Soll} &
  \headercol{1in}{Ist} \\
 \end{tabular}
 \caption{Aufw�nde}
 \label{table:aufw�nde}
\end{table}

Eine genauere Analyse der Aufw�nde findet man im Anhang Statistik (siehe \ref{app:statistik}).


\subsection{Kosten}
Die geplanten Kosten entsprechen in etwa den vorgesehenen Projektkosten, da der gr�sste Kostenpunkt, die variablen Kosten, ziemlich genau der Planung entspricht.



\section{Erkenntnisse und Erfahrungen}
\subsection{Einarbeiten in das Themengebiet}
\subsection{Evaluation Algorithmen}
\subsection{Implementation Algorithmus}
\subsection{Implementation Testframework}
\subsection{Analyse GUI/Netzwerk}

\subsection{Hermes}


\subsection{Coaching Projektf�hrung}
Die M�glichkeit des Projektf�hrungs-Coaching durch Herrn Helbling haben wir rege in Anspruch genommen. Diverse Fragen stellten wir auch per e-mail. Allgemein konnten wir einige wichtige Erkenntnisse daraus gewinnen. 

Herr Helbling hat uns zu Beginn darauf aufmerksam gemacht, dass er einen gewissen Interessenkonflikt zwischen seiner Rolle als PM-Coach und seiner Rolle als Notengeber hat. Wir haben die Auswirkungen dieses Interessenkonflikt selber erlebt. Gewisse Fragen haben wir mehrmals gestellt und nie eine wirklich befriedigende Antwort erhalten. Verweise auf Hermes-Abschnitte sind in der Regel nicht sehr hilfreich, da man die betreffenden Punkte selber schon in Hermes angeschaut hat. 

Unserer Meinung sollte man, wenn dieser Interessenkonflikt wirklich so gross ist, eine �nderung am System vornehmen. Es sollte evtl. ein anderer Dozent PM-Coach sein als der, der die Noten verteilt. Dies w�rde ein unvoreingenommenes Coaching erm�glichen, was sicher f�r alle beteiligten befriedigender w�re.


\subsection{Coaching durch Fachdozenten}
Die Betreuung durch den Fachdozenten wurde w�hrend der Projektarbeit nur selten ben�tigt. Bei Fragen standen uns Herr Dubois und Herr Fuhrer zur Verf�gung, was die Zusammenarbeit erleichterte und ein z�giges Arbeiten erm�glichte. 


\subsection{Projektarbeit}



\section{Pers�nliche Bilanzen}
\subsection{Mark Bigler}
Ich bin sehr froh dar�ber, dass unser Vorschlag f�r das Semesterprojekt angenommen wurde. Dadurch konnten wir auf einem Gebiet arbeiten, welches jeden von uns interessiert hatte. Zus�tzlich haben alle von uns mehr Schulstunden im Wintersemester als im Sommersemester ausgew�hlt, was uns die M�glichkeit gab, die Zeit f�r Arbeiten am Projekt relativ flexibel zu gestallten. Im grossen und ganzen war es eine sehr gute Erfahrung. Zum Beispiel hatte ich w�hrend der Implementation des Algorithmus zum ersten Mal die M�glichkeit gelernte Software-Patterns aus dem zweiten Studienjahr anzuwenden.

\subsubsection{Positive Punkte}
\begin{itemize}
 \item Wir haben uns ein interessantes Thema ausgew�hlt und waren somit sehr motiviert.
 \item Dank dem Zeitplan konnten wir die Arbeit gut ein- und aufteilen.
 \item Wir konnte uns realistische Ziele setzen und haben diese auch alle erreicht.
\end{itemize}

\subsubsection{Negative Punkte}
\begin{itemize}
 \item Sehr viel Aufwand f�r Projekt Management Dokumente. Wahrscheinlich h�tten wir sehr gut weniger machen k�nnen, ohne Konsequenzen f�r den Verlauf oder die Qualit�t des Projektes.
\end{itemize}

\subsubsection{Selbstkritik}
\begin{itemize}
 \item Ich habe die Zeit w�hrend den Fr�hlingsferien nicht optimal ausgenutzt. Deshalb konnte ich den GUI Bericht auch nicht bis zum geplanten Termin fertigstellen. Dies hatte aber keinen Einfluss auf weiteren Aktivit�ten des Projektes.
\end{itemize}
\subsection{Simon R�ss}
Ein Projekt ist immer eine Bereicherung und Abwechslung zum Unterricht. Gl�cklicherweise hatten wir im Sommersemester deutlich weniger Stunden als von der Schule vorgesehen und konnten uns dadurch im Allgemeinen voll auf das Projekt konzentrieren. Das hatte den positiven Effekt, dass wir alle unseren gesetzten Ziele, sogar die Ziele mit Priorit�t 2, erreichen konnten. Im Hinblick auf die Diplomarbeit kann sicher noch die Teamzusammenarbeit verbessert werden, zum Beispiel durch regelm�ssigere Meetings.

\subsubsection{Positive Punkte}
\begin{itemize}
 \item Freiheit beim Festlegen der Ziele und der Realisierung
 \item Sehr interessantes und innovatives Themengebiet
 \item Alle gesetzten Ziele wurden erreicht
\end{itemize}

\subsubsection{Negative Punkte}
\begin{itemize}
 \item Aufwand f�r Projektmanagement Dokumente zu hoch
 \item Ich hatte mehr Nutzen erwartet vom Projektmanagement Coaching.
\end{itemize}

\subsubsection{Selbstkritik}
Die Effizienz der Zusammenarbeit k�nnte sicher durch regelm�ssige Teammeetings noch verbessert werden. Insbesondere als Teamleiter muss man wissen, wie der Fortschritt des Projektes ist und wo noch Probleme bestehen, damit man fr�h Korrekturmassnahmen einleiten kann.

\subsection{Lukas Zbinden}




\section{Fazit}
Alle unseren gesteckten Ziele wurden erreicht. Damit sind wir bereit, in der Diplomarbeit das eigentliche Ziel von \ace, einen kollaborativen Texteditor zu entwickeln, anzugehen.

Es kann auch gesagt werden, dass das Projekt Spass gemacht hat und wir dabei viel lernen konnten. Dank der Projektarbeit konnten wir wichtige praktische Erfahrungen sammeln und uns in ein interessantes Themengebiet einarbeiten.

F�r das gesamte \ace Projektteam steht fest, dass wir das Projekt als Diplomarbeit fortsetzen werden.



\newpage
\appendix

\section{Undo in multiuser collaborative applications}
\label{sect:undo}

The availability of undo in a multiuser collaborative application is valuable because features available in single-user applications should also be available in corresponding multi-user applications. However, supporting undo in a multiuser collaborative application is much more difficult than supporting undo in single-user interactive applications because of the interleaving of operations performed by multiple users in a collaborative computing environment.


\subsection{An Example}

To illustrate the difficulties, let us consider a collaborative editing session with two users and a shared text document containing the string '"abc'". Suppose user 1 issues the operation $O_{1} = Ins[2,X]$ to insert the character '"X'" at position 2 after the character '"a'". The resulting document state is '"aXbc'". After this, user 2 issues operation $O_{2} = Ins[2,Y]$ to insert character '"Y'" after the character '"a'" to get the resulting document state '"aYXbc'". Suppose user 1 wants to undo his last operation, i.e. $O_{1}$. Her intended effect is to remove the '"X'" thus resulting in a document state '"aYbc'". First, blindly picking the \emph{globally} last operation, i.e. $O_{2}$ will undo the wrong operation. Second, simply executing the inverse operation $\overline{O_{1}} = Del[2,X]$ to delete the '"X'" at position 2 in the current document state '"aYXbc'" will delete '"Y'" instead of '"X'".

This means that any group undo solution must meet the challenge of undoing operations in a nonlinear way and must be able to achieve the correct undo effect in a document state that has been changed by other users'' operations. It can be clearly seen that undo in a groupware system is highly related to operational transformations.


\subsection{Types of Undo}

Generally a user expects an undo to reverse their own last operation (\emph{local} undo) rather than the globally last operation (\emph{global} undo). So an undo framework for groupware systems needs to allow selection of an operation to undo based on who performed it. Undoing the globally last action can be problematic as just before one user presses undo, another user may issue an operation. The effect of the operation may get executed at the first users site just before he effectively invokes the undo action. This would not result in the undoing of an operation which the first user intended to undo.

Further an undo mechanism may be classified by whether it allows to undo an arbitrary operation (called \emph{selective} undo) or it is restricted to the \emph{chronological} order.

It is still an open question how to select operations to be undone in a selective undo system. To undo the chronoligically last operation of the local user, the single familiar undo menu or shortcut is sufficient as it is always clear wich operation has to be undone.


\subsection{Existing Undo Solutions}

\subsubsection{DistEdit System}
The DistEdit system \cite{prakash94} allows operations to be undone in any order. However, it may be that an operation is not undoable because it conflicts with a later executed operations. The conflicts occur if a later operation happens at the same index as the operation to undo. They managed to solve some conflicting situations, but still a few remained in the finished system.

\subsubsection{adOPTed undo}
In the \emph{adOPTed} algorithm (see \ref{algo:adopted}) undo is supported by operational transformation. It allows to undo the chronological last operation of a user (it is implemented only for the local user but would it would be possible to extend it to an arbitrary user). To undo an operation, an inverse of this operation is generated by a special \emph{mirror} operator. This inverse operation must be placed at a valid location in the corresponding dimension of the interaction graph. No conflict can occur in undoing any operation. Another special operator called \emph{folding} is used for a correct working system. See \cite{ressel99} for a detailed description of this undo system.

\subsubsection{GOTO undo}
Sun \cite{sun02b} described an undo mechanism for the \emph{REDUCE} prototype. It is using the \emph{GOTO} algorithm (see \ref{algo:goto}) for \emph{do} and an algorithm called \emph{ANYUNDO} for \emph{undo}. The algorithm allows undoing arbitrary operations, so it is a selective undo system. The algorithm is described in great detail along many possible problems and how they were solved in this algorithm.



\section{Vector Time}
\label{sect:vectortime}

Most algorithms use \emph{vector time} to determine causal relations. Each site maintains a state vector $v$ that has $n$ components. $n$ is the number of participating sites. $S$ is the local site. The $i$-th component of $v$ denoted as $v[i]$ represents the number of operations executed from site $i$ at site $S$ (the local site). 

\begin{defn}
  For two time vectors u, v \\
  $u \leq v$ iff $\forall i : u[i] \leq v[i]$ \\
  $u < v$ iff $u \leq v$ and $u \not= v$ \\
  $u \parallel v$ iff $\neg(u < v)$ and $\neg(v < u)$
\end{defn}

Notice that $\leq$ and $<$ are partial orders. The \emph{concurrency relation} $\parallel$ is reflexive and symmetric.

The above definition gives us a very simple method to decide whether two events $e$ and $e'$ are causally related or not: We take their timestamps (vector times) and check whether $C(e) < C(e')$ or $C(e') < C(e)$ where $C(x)$ determines the timestamp of $x$. If the test succeeds, the events are causally related. Otherwise they are independent.




\end{document}
